\documentclass[11pt]{article}
\usepackage[margin=1in]{geometry}
\usepackage{amsmath,amssymb,amsthm}
\usepackage{booktabs}
\usepackage[hidelinks]{hyperref}

\newtheorem{theorem}{Theorem}
\newcommand{\cE}{\mathcal{E}}
\newcommand{\cB}{\mathcal{B}}
\newcommand{\cA}{\mathcal{A}}
\newcommand{\cK}{\mathcal{K}}
\newcommand{\cF}{\mathcal{F}}
\newcommand{\cQ}{\mathcal{Q}}
\newcommand{\cT}{\mathcal{T}}
\newcommand{\sym}{\operatorname{sym}}
\newcommand{\Area}{\operatorname{Area}}
\newcommand{\Vol}{\operatorname{Vol}}
\newcommand{\Srel}{S_{\rm rel}}
\newcommand{\GD}{G_{\!D}}
\newcommand{\diam}{\Diamond_B}
\newcommand{\ip}[2]{\langle #1,#2\rangle}
% evidence tag, defined in Section 1
\newcommand{\ev}[1]{\,{\textsf{\scriptsize[#1]}}}
% ragged-right table cell without the array package
\newcommand{\rr}[1]{\parbox[t]{\linewidth}{\raggedright #1}}
% in-text tables (kept in place, not floated)
\newcommand{\tabcap}[2]{\par\smallskip\refstepcounter{table}\label{#1}{\small Table~\thetable: #2}\par}
% tighter section spacing than the article default
\makeatletter
\renewcommand\section{\@startsection{section}{1}{\z@}{-2.4ex \@plus -1ex \@minus -.2ex}{1.2ex \@plus .2ex}{\normalfont\Large\bfseries}}
\makeatother

\title{The second-order area charge of a flat ball generates no normal\\ modular flow on the linearised-curvature algebra}
\author{Andreas Houg}
\date{16 November 2026 (version 4)}

\begin{document}
\maketitle

\begin{abstract}
We take the observables of free linearised gravity in a flat ball to be the algebra $M$ generated by the linearised curvature, and we take the second-order area charge to be the second-order part of $\Area-2\Vol$ at its saddle over the edge's shape, the corner's time shift and the slice. On each multipole sector this charge equals the vacuum's modular generator plus a positive form $Z$ on extrinsic-curvature data, local for odd parity and a Dirichlet Green form for even parity; we prove this for every degree $l\ge2$, given one inherited completeness statement for the transverse-traceless data. For every coupling $g\neq0$ and every modular time $t\neq0$, the flow of the vacuum generator plus $g$ times $Z$ is then implemented by no unitary on Fock space, because the Hilbert--Schmidt norm of its antilinear part has a divergent harmonic tail over sectors. Hence no faithful normal weight or state on $M$ has the area-side flow ($g=1$) as its modular group at any non-zero time. Every lane, check and referee behind these statements is an AI model of one family, working under the author's direction; nothing here is peer reviewed.
\end{abstract}

\section{Result, assumptions and evidence}\label{sec:result}

Units: $\hbar=c=1$, $G$ kept. $B$ is the unit ball in the slice $t=0$ of Minkowski space (radius $R$ by dilation), $\diam$ its causal diamond, and the theory is free linearised gravity at order $G^0$ in four dimensions. A finite-norm solution defines a coherent state; $\Srel$ is its relative entropy to the vacuum on the observables of $\diam$. Each multipole sector has a degree $l\ge2$ and a parity, even (polar) or odd (axial); its Cauchy data split into \emph{metric data} $h_{ij}$ and \emph{extrinsic-curvature data} $K_{ij}$. With $\cE_{ij}=R_{0i0j}$ and $\cB_{ij}=\tfrac12\epsilon_{ikl}R_{0j}{}^{kl}$, the sector's \emph{chain field} is $\psi=x^ix^j\cE_{ij}$ (even) or $x^ix^j\cB_{ij}$ (odd), of the form $f(r)Y_{lm}$; for even parity $\psi$ is built from metric data and $\dot\psi$ from extrinsic curvature, for odd parity the other way round. With $N=\frac12(1-r^2)$, $c_l=4(l-2)!/(l+2)!$ and $\GD$ the Dirichlet Green operator of $-\nabla^2$ on $B$, define
\[
E_w[\psi,\dot\psi]=\tfrac12\int_BN\bigl(\dot\psi^2+|\nabla\psi|^2\bigr)+\tfrac w2\int_B\psi^2 .
\]

\paragraph{Tags.} Each claim ends in a bracket saying how it was checked. \textsf{\small P}: proved in the record (a written proof whose core the programme's lead checked by hand, with a script re-run); \textsf{\small X}: exact arithmetic; \textsf{\small N}: certified or high-precision numerics; \textsf{\small L}: lattice; \textsf{\small S1}, \textsf{\small S2}, \textsf{\small S3}: a standard theorem read in a primary source, in a secondary source, or quoted from memory; \textsf{\small E}: carried from the earlier record, not re-examined; \textsf{\small Lit}: a reading of the literature; \textsf{\small ref}: a hostile referee found no break. The word after the colon names the report in the record~\cite{T}. Every lane, the lead's checks, the referees and the second codes are AI models of one family working under the author's direction.

\paragraph{The result.} Let $\cA$ be the second-order area charge of Section~\ref{sec:charge} (in units of $1/8G$). On every sector, both parities and every $l\ge2$,
\begin{equation}\label{eq:sector}
\text{odd: }\ \cA=c_lE_2=c_lE_1+\tfrac12c_l\!\int_B\psi^2,\qquad
\text{even: }\ \cA=c_lE_1+\tfrac12c_l\,\ip{\dot\psi}{\GD\dot\psi},
\end{equation}
so $\cA=c_lE_1+Z$ with $Z\ge0$ vanishing on metric data and positive on extrinsic-curvature data.\ev{P: charge-all-l, symbolic in $l$; ref: charge-ref; given I} Given the definition A1 below, $2\pi c_lE_1$ is the vacuum's modular generator on the sector (Section~\ref{sec:setting}). Theorem~\ref{thm:main} (Section~\ref{sec:shale}) shows that for every real $g\neq0$ and every $t\neq0$ the flow of $c_lE_1+gZ$ over the whole tower is implemented by no unitary on Fock space, so it is the modular flow of no faithful normal semifinite weight or state on $M$ at time $t$. If A1 and A2 hold, the area-side flow ($g=1$) is therefore the modular flow of no faithful normal state or weight of the ball's graviton algebra in the vacuum sector. ``Normal'' is the scope: states of the Weyl C$^*$-algebra outside the vacuum representation are not excluded, and one exists (Section~\ref{sec:shale}).

\paragraph{What it rests on.} Two choices are assumed: \textbf{A1}, the ball's graviton algebra is $M$, the algebra generated by the linearised curvature on the transverse-traceless Fock space; \textbf{A2}, this saddle value of the second-order $\Area-2\Vol$ is the physically relevant charge. One step is inherited: \textbf{I}, the two families of transverse-traceless sector data used in Section~\ref{sec:charge} exhaust the smooth data of a sector. Standard theorems are sorted in Table~\ref{tab:hyp}; the cone paragraph of Section~\ref{sec:ruler} adds two hypotheses of its own.

\medskip\noindent\begin{minipage}{\textwidth}\centering\footnotesize
\tabcap{tab:hyp}{What the results rest on. ``Proved'' means a written proof whose core was checked by hand by the programme's lead, with a script re-run; the reports are listed in~\cite{T}.}\smallskip
\begin{tabular}{p{0.16\textwidth}p{0.79\textwidth}}
\toprule
Status & Items\\
\midrule
\rr{\textbf{Proved}} & The sector identity~\eqref{eq:sector} for every $l$, as exact radial identities symbolic in $l$, given I (refereed inside the programme). The change of variables to the hyperbolic frame, coupling exactly~1, both parities. $\|D\|\le\sqrt2/\lambda$. Density of smooth data, boundedness of $D$, isometric vacuum flow, $\operatorname{Ker}h=\{0\}$. Orientation of the vacuum modular flow (Tomita computation). The antilinear part, exactly. An implementer forces a bounded symplectic extension. First-order floors on explicit families, uniform in $l\ge16$. Second-order remainders. Every $t\neq0$. The tower sum. The constants recomputed by blind second codes with different methods.\\
\rr{\textbf{Graviton bridge, proved}} & Given A1: the vacuum modular group of the double cone acts on curvature smeared inside it as the conformal flow, weight~3 (refereed). For every $l$, each sector's curvature is determined by its chain field, a massless scalar, with the inner product on data $c_l$ times the scalar's (refereed). With Hislop--Longo and second quantisation: the sector generator $2\pi c_lE_1$.\\
\rr{\textbf{Standard, primary source read}} & Hislop--Longo~\cite{HL} Thm~2, massless scalar. Connes~\cite{Connes} Thm~1.2.1, Lemma~1.2.2(c), forward direction only. Brunetti--Guido--Longo~\cite{BGL} Thm~2.3(ii). Lled\'o~\cite{Lle} Thm~5.7, Prop.~5.8. Longo--Morsella~\cite{LM}. (One reading each, not repeated.)\\
\rr{\textbf{Standard, secondary source}} & Shale, necessity direction~\cite{Shale} (Dere\-zi\'nski~\cite{Der} Thms~23--24). Modular operator of a second-quantised algebra (Dere\-zi\'nski Thm~40(6), needs $\operatorname{Ker}h=\{0\}$). Bounded perturbation (Schnaubelt~\cite{Schn} Thm~3.4). Totality of coherent vectors. L\"owner--Heinz.\\
\rr{\textbf{Standard, from memory}} & Limit-point criterion (Reed--Simon X.10), for self-adjointness of $h^2$. Weyl's criterion and Reed--Simon X.7, used for density. The factoriality fact in the vacuum covariance. The time-slice step (data on the ball span the same one-particle subspace as test functions in the double cone). Type III of free-field local algebras (Section~\ref{sec:lit} only).\\
\rr{\textbf{Inherited}} & I above (completeness of the two data families).\\
\rr{\textbf{Assumed}} & A1 (a definition) and A2 (a physical choice).\\
\rr{\textbf{Open}} & Section~\ref{sec:open}.\\
\bottomrule
\end{tabular}
\end{minipage}

\section{Setting: flows and the vacuum generator}\label{sec:setting}

Let $\zeta=\frac12(1-t^2-r^2)\partial_t-t\,x\!\cdot\!\partial$ be the conformal Killing vector of $\diam$; its lapse on the slice is $N$, and $\nabla\!\cdot\zeta=-4t$. On the chain field, the \emph{weight-$w$ flow} is $\delta\psi=\zeta\psi-wt\psi$; on the slice
\begin{equation}\label{eq:flow}
\delta\psi=N\dot\psi,\qquad \delta\dot\psi=\nabla\!\cdot(N\nabla\psi)-w\,\psi ,
\end{equation}
and $E_w$ generates it. Three flows appear below; we name them once.
\begin{itemize}\itemsep1pt
\item \textbf{F1, the vacuum modular flow.} Given A1, $\Delta^{it}$ acts on curvature smeared inside the double cone as the conformal transport of the curvature by $\zeta$ with conformal weight~3, at flow parameter $-2\pi t$.\ev{P: wave70 modflow; ref: modflow-ref} On chain data this is the weight-1 flow~\eqref{eq:flow}, generated by $c_lE_1$.\ev{P: modgen, chain-all-l}
\item \textbf{F2, the area-side flow,} generated by $\cA$: the weight-2 flow for odd parity; for even parity the weight-1 flow plus the flow of $\frac12c_l\ip{\dot\psi}{\GD\dot\psi}$.
\item \textbf{F3, Lie transport of the metric perturbation,} $\delta h_{ij}=\pounds_\zeta h_{ij}+2t\,h_{ij}$ (the background rescaling removed). On odd sectors its gauge-invariant canonical charge is $c_lE_2$, the same as F2.\ev{X: charge-id, jss-id, $l=2$--$6$}
\end{itemize}
Which weight a ``diffeomorphism acting on the curvature'' carries depends on how the background's conformal factor is treated; F1 is the transport that is a symmetry of the free curvature field.

\paragraph{The bridge.} For every $l$ the chain field obeys the massless scalar wave equation and determines the sector's curvature and is determined by it; $\psi$ and $\dot\psi$ smeared in the ball generate the same one-particle local subspace as the sector's curvature; the map from curvature data to chain data is symplectic and isometric up to $c_l$ and intertwines time evolution.\ev{P: chain-all-l; ref: chain-ref} The constant multiplies the inner product on data; the chain field's own two-point function is $1/c_l$ times the scalar's. A transverse-traceless Cauchy datum is fixed by $\psi$ only up to residual gauge modes; the curvature is fixed exactly. With Hislop--Longo~\cite{HL} (unit cone, radius $R$ by dilation, plus the time-slice step) and second quantisation, the vacuum's modular generator on each sector is $2\pi c_lE_1$, and for a coherent state $8G\,\Srel=\sum_{\rm sectors}c_lE_1[\psi,\dot\psi]$.\ev{P: modgen, subspace; S1, S2, S3; given A1}

\paragraph{Sectors and parity.} Rotations and parity commute with $E_1$, $\cA$ and the vacuum, and the (parity, $l$) data carry inequivalent irreducible representations of $O(3)$, so by Schur's lemma all these forms are block diagonal in (parity, $l$), act as the identity on $m$, and have no cross terms between parities or between different $l$. We write $M_\pm$ for the algebra generated by the curvature of one parity; $M=M_+\vee M_-$, and since the maps below are block diagonal, non-implementability for one parity gives it for $M$.

\paragraph{Hyperbolic frame.} Put $r=\tanh(\rho/2)$, so $N=1/(1+\cosh\rho)$ and $dr=N\,d\rho$; $q=rf$, $\pi=Nr\dot f$, $\lambda^2=l(l+1)$, $h=h_l$ with $h_l^2=-\partial_\rho^2+\lambda^2/\sinh^2\!\rho$ on $L^2(d\rho)$. Then $E_1=\tfrac12\ip\pi\pi+\tfrac12\ip q{h^2q}$, and~\eqref{eq:flow} at $w=1$ is exactly $2\pi X_0$ per unit of modular time, $X_0(q,\pi)=(\pi,-h^2q)$; the vacuum covariance is the thermal form of $h$ at inverse temperature $2\pi$.\ev{P: modgen; N: 23 digits, $l=2,3$} In these variables $Z$ has the Hamiltonian vector field $D_{\rm even}(q,\pi)=(\GD\pi,0)$, $D_{\rm odd}(q,\pi)=(0,-Nq)$, and $\cA=X_0+D$ with coupling exactly~1, for every $l$.\ev{P: even-dict, odd-dict}

\section{The area-side charge}\label{sec:charge}

For a small change $q_{ij}$ of the flat metric, Green's identity with $\nabla\nabla N=-\delta_{ij}$ and $N|_{\partial B}=0$ gives $\delta\Area-2\delta\Vol=-\frac12\int_BN\,R_{\rm lin}[q]$; with the Hamiltonian constraint this is the first law of causal diamonds~\cite{JV}, and at second order it removes the second-order metric.

\paragraph{Definition.} Move the edge by a normal shift $\sigma$, shift the corner in time by $\theta$ and choose the slice by a function $X$. Each choice is a gauge transformation of the data, and fixed true volume equals $\Area-2\Vol$ at second order\ev{P: charge-ref}; the unextremised value $\cA_{\rm un}$ ($-2$ times the second-order part of $\Area-2\Vol$, in units of $1/8G$) depends on the choice, and its stationary value over $(\sigma,\theta,X)$ is the one number that does not. That value is $\cA$. Stationarity in $X$ makes the slice maximal, and the stationary point is a saddle: a maximum in the edge shape and the slice, a minimum in the time shift. Without the time shift the even value is gauge dependent; including it is our reading of Jensen--Sorce--Speranza's Section~3.3~\cite{JSS23}.\ev{X: even-tower, $l=2$--$6$}

\paragraph{Proof of~\eqref{eq:sector}, in outline.} Integrals are per $\oint Y^2$, with $\varphi=r^lY$.
\begin{enumerate}\itemsep0pt\setlength{\leftmargin}{1.6em}
\item[(i)] On transverse-traceless data, and only there (confirmed term by term by the referee),
\[\textstyle\cA_{\rm un}=\int_BN\,K_{ij}K_{ij}+\frac14\int_BN|\nabla h_{ij}|^2-\frac12\oint|h_{ij}\nu_j|^2+\frac14\oint(\nu_ih_{ij}\nu_j)^2.\]
\item[(ii)] \emph{Edge.} The first-order area change has the drive $\oint\sigma G$ with $G=\frac12(\partial_r\psi_h+2\psi_h)|_{r=1}$, $\psi_h=x_ih_{ij}x_j$, and the stiffness is $\frac12(l-1)(l+2)\oint\sigma^2$, so the edge adds $\oint G^2/((l-1)(l+2))$.
\item[(iii)] \emph{Corner.} For transverse-traceless $K$ and any $X$, $\int_BN(|K+\nabla\nabla X|^2-(\nabla^2X)^2)=\int_BN|K|^2+2\oint\kappa\theta+\oint|\nabla_{\!\rm tan}\theta|^2-2\int_B|\nabla X|^2$, with $\theta=X|_{\partial B}$ and $\kappa=x\!\cdot\!K\!\cdot\!x$. At the stationary point $X$ is harmonic, $\theta=-\kappa/(l(l-1))$, and the value is $\int_BN|K|^2-\oint\kappa^2/(l(l-1))$.
\item[(iv)] \emph{Even, extrinsic block.} Write $K=T[q]$, the transverse-traceless even tensor with $x\!\cdot\!T\!\cdot\!x=q\varphi$, and $\tilde q=q''+(2l+2)q'/r$, so $\dot\psi=-\tilde q\varphi$. Pointwise in $r$, in the record's radial normalisation, $|T[q]|^2_r-\frac{c_l}2\tilde q^2=r^{-2l-2}(r^{2l+2}U)'$ with $U=\bigl((2l+1)q^2/r^3+2qq'/r^2+2q'^2/(r(l+1)(l+2))\bigr)/(l(l-1))$. Integrating against $N$ ($N'=-r$) gives $\int_BN|K|^2=\frac{c_l}2\int_BN\dot\psi^2+q(1)^2/(l(l-1))+\frac{c_l}2\int_0^1r^{2l+2}q'^2dr$. The corner term cancels $q(1)^2/(l(l-1))$ exactly, and since $\nabla^2((q-q(1))\varphi)=\tilde q\varphi$, $\GD\dot\psi=(q-q(1))\varphi$ and the last term is $\frac12c_l\ip{\dot\psi}{\GD\dot\psi}$.
\item[(v)] \emph{Even, metric block.} The density of $\cA_{\rm un}$ minus that of $c_lE_1$ is a total derivative $r^{-2l-2}(r^{2l+3}W)'$ with a potential $W$ rational in $l$, whose boundary value the edge term cancels; the stationary edge shape on the static tide is $2/(l+2)$.
\item[(vi)] \emph{Odd.} Odd data have parity $(-1)^{l+1}$ and no scalar of degree $l$ has it, so there is no drive and $\cA=\cA_{\rm un}$; the same kind of total-derivative identity gives $c_lE_2$, with $\psi$ built from $K$ alone.
\item[(vii)] \emph{Cross block.} Under $t\to-t$, $h_{ij}$ drives only $\sigma$ and $K$ only $(\theta,X)$, so there is no metric--extrinsic cross term.\ev{P: charge-all-l, symbolic in $l$; X: Cartesian, $l=2$--$9$; ref: charge-ref; given I}
\end{enumerate}

\paragraph{The ladder.} Against $c_lE_1$ on extrinsic-curvature data, $\cA$ has the eigenvalues
\begin{equation}\label{eq:ladder}
1+\frac{2}{n(2l+n+1)},\qquad n\ \text{odd for odd parity, even for even parity}.
\end{equation}
For odd parity the eigenproblem is $(L+2)\psi=\mu(L+1)\psi$ with $L=-\nabla\!\cdot(N\nabla\,\cdot)$. Since $\nabla\!\cdot(N\nabla(r^kY))=\frac12a_kr^{k-2}Y-(\frac12a_k+k)r^kY$ with $a_k=(k-l)(k+l+1)$, $L$ is triangular on $r^{l+2j}Y$ with diagonal $\nu_j=j(2l+2j+1)+l+2j$, and $\mu=1+1/(\nu_j+1)$ is~\eqref{eq:ladder} with $n=2j+1$. For even parity the eigenproblem is $-N\nabla^2u=u/(\mu-1)$ with $u=\GD\dot\psi$; $-N\nabla^2$ is triangular on $Nr^{l+2j}Y$ with diagonal $j(2l+2j+1)+2l+4j+3$, which is~\eqref{eq:ladder} with $n=2j+2$. The eigenvectors span the polynomial sector data, which are dense. Rung $n=1$ is $(l+2)/(l+1)$.\ev{derived here from~\eqref{eq:sector}; X: even-tower, charge-all-l, first three rungs}

\paragraph{The symptom.} Version~3 compared two tides. Here $t$ is Minkowski time. The static electric tide ($\Phi$ harmonic of degree $l$, $ds^2=-(1+2\Phi)dt^2+(1-2\Phi)d\mathbf x^2$, $\psi=l(l-1)\Phi$) is pure even metric data; the stationary magnetic tide ($h_{ij}=2\varepsilon t\,\sym\nabla(x\times\nabla\phi)_{ij}$, $\psi=\frac12(l-1)l(l+1)\varepsilon\phi$) is pure odd extrinsic-curvature data. For static data $E_w=\frac{l+w}2\int_B\psi^2$, so
\begin{equation}\label{eq:S}
8G\,\Srel=c_l\tfrac{l+1}2\textstyle\int_B(\psi_\cE^2+\psi_\cB^2),\qquad
8G\,(\text{area side})=c_l\tfrac{l+1}2\textstyle\int_B\psi_\cE^2\ \text{or}\ c_l\tfrac{l+2}2\textstyle\int_B\psi_\cB^2 .
\end{equation}
The static electric tide is pure metric data, so $Z$ vanishes on it and the two sides agree when the edge sits at its stationary shape; on a geodesic ball (edge not at that shape) the ratio at $l=2$ is $2/9$.\ev{E; X: ruler-why} On stationary magnetic data, and only there, entropy/area side $=(l+1)/(l+2)$, the inverse of rung $n=1$. For $l=2$, $\int_B\psi_\cE^2=\frac{8\pi}{105}\cE_{ij}\cE_{ij}$ gives $\Srel=\frac{\pi R^6}{420G}(\cE_{ij}\cE_{ij}+\cB_{ij}\cB_{ij})$, against the area terms $\frac{\pi R^6}{420G}\cE_{ij}\cE_{ij}$ and $\frac43\frac{\pi R^6}{420G}\cB_{ij}\cB_{ij}$, which are in print~\cite{JSS17}.\ev{E}

\section{Non-implementability}\label{sec:shale}

$G$ and $c_l$ rescale each sector's data by a constant and cannot change whether an operator is Hilbert--Schmidt; this section omits them and fixes one parity. From here on $t$ is modular time.

\paragraph{Setting.} Per sector let $L=L^2(d\rho)$ and $z=(h^{1/2}q+ih^{-1/2}\pi)/\sqrt2$. Define first $\cK=\bigoplus_{l,m}(L\oplus\bar L)$ and in it
\[
k(z)=(c_hz,\,s_h\bar z),\quad k'(w)=(s_hw,\,c_h\bar w),\qquad c_h=(1-e^{-2\pi h})^{-1/2},\ s_h=(e^{2\pi h}-1)^{-1/2},
\]
so $c_h^2-s_h^2=1$, $c_h^2+s_h^2=\coth\pi h$, $c_h/s_h=e^{\pi h}$. $H$ is the closed span of the $k(z)$, with the thermal norm $\|k\|^2=\ip z{\coth(\pi h)z}$ and $\operatorname{Im}\ip{k(z)}{k(z')}=\frac12\omega$, $\omega=\ip q{\pi'}-\ip\pi{q'}$; $H'$, spanned by the $k'(w)$, is its symplectic complement. $H+H'$ is dense in $\cK$: every $(u,v)$ in the dense domain of $c_h$ equals $k(c_hu-s_h\bar v)+k'(c_h\bar v-s_hu)$.\ev{P: beta-def} On $\cF(\cK)$, $W(k)=e^{i(a(k)+a^*(k))}$, $\ip\Omega{W(k)\Omega}=e^{-\|k\|^2/2}$ (Dere\-zi\'nski's $k$ is $\sqrt2$ times ours), $M_\pm=\{W(k):k\in H\}''$, and $\omega_0$ is the vacuum state. $\operatorname{Ker}h=\{0\}$, since $\ip f{h^2f}=\|f'\|^2+\lambda^2\|f/\sinh\rho\|^2$ vanishes only for $f=0$; smooth compactly supported data are dense in $H$, $D$ is bounded, and $U^0(s)=e^{sX_0}$, which acts as $z\mapsto e^{-ish}z$, is a strongly continuous isometry group.\ev{P: even-density; S3} $\|D\|\le\sqrt2/\lambda$.\ev{P: even-proof, hinge-proof} With $U^g(s)=e^{s(X_0+gD)}$ and $\tau=2\pi t$,
\[
S(\tau)=U^g(\tau)\,U^0(-\tau),\qquad \hat S=S\oplus1\ \text{on}\ H+H',\qquad \beta=\tfrac12(\hat S+i\hat Si);
\]
writing $(S-1)z=Az+B\bar z$, exactly $\beta k(z)=(c_hP,s_h\bar Q)$ with $P=B(c_h^2\bar z)-A(s_h^2z)$, $Q=A(c_h^2z)-B(s_h^2\bar z)$.\ev{P: beta-def}

\begin{theorem}\label{thm:main}
For either parity, every real $g\neq0$ and every $t\neq0$, $\hat S(2\pi t)$ is implemented by no unitary on $\cF(\cK)$. Hence no faithful normal semifinite weight on $M_\pm$ (or on $M$), in particular no faithful normal state, has modular automorphism $\sigma_t(W(k))=W(U^g(2\pi t)k)$ for all $k\in H$. Quantitatively, there are unit vectors $\hat e_{l,m}$, orthonormal over $\mathbb C$, with $\|\beta\hat e_{l,m}\|\ge c(\tau)|g|/(2\lambda)$ for all $l\ge l_1(g,\tau)$, where
\[
c(\tau)=c_\pi\,\Upsilon\bigl(\min(|\tau|,s_*)/s_*\bigr),\quad \Upsilon(\sigma)=\tfrac18(15\sigma-10\sigma^3+3\sigma^5),\quad s_*=1.777101,
\]
\[
l_1(g,\tau)=\max\bigl(16,\ \lceil 2e^{\eta}C_R\,\tau^2|g|/c(\tau)\rceil\bigr),\qquad \eta=15c_\pi/(8\sqrt2\,s_*C_R)\approx10^{-3},
\]
with $c_\pi$ and $C_R$ in Table~\ref{tab:const}.\ev{P: assembly, referee, all-t and steps (a)--(f); S1: Connes; S2: Shale, Dere\-zi\'nski Thm 40(6), bounded perturbation, totality; S3: limit point, Weyl}
\end{theorem}

\paragraph{(a) Vacuum modular flow.} $H\cap iH=0$ and $H+iH$ is dense. The Tomita operator $\cT_H:k(z)+ik(z')\mapsto k(z)-ik(z')$: writing $k(z)+ik(z')=(u,v)$, $u=c_h(z+iz')$, $v=s_h\overline{(z-iz')}$, one gets $\cT_H(u,v)=((c_h/s_h)\bar v,(s_h/c_h)\bar u)$, with adjoint $(a,b)\mapsto((s_h/c_h)\bar b,(c_h/s_h)\bar a)$, so
\[
\Delta_H=\cT_H^*\cT_H=e^{-2\pi h}\oplus e^{2\pi h},\qquad \Delta_H^{it}k(z)=k(e^{-2\pi ith}z)=k\bigl(U^0(2\pi t)z\bigr).
\]
Dere\-zi\'nski's Theorem~40(6), with $\gamma=e^{-2\pi h}$ and $\operatorname{Ker}h=\{0\}$, gives $\Delta=\Gamma(\Delta_H)$, so $\sigma^{\omega_0}_t(W(k))=W(U^0(2\pi t)k)$.\ev{P: referee; S2}

\paragraph{(b) Weight $\Rightarrow$ implementer.} If a faithful normal semifinite weight $\varphi$ had $\sigma^\varphi_t$ equal to $W(k)\mapsto W(U^g(\tau)k)$ at one $t$, Connes' theorem gives a unitary $u\in M_\pm$ such that $\sigma^\varphi_t\circ\sigma^{\omega_0}_{-t}$ is $\operatorname{Ad}u$, which on $W(k)$ is $W(S(\tau)k)$; as $W(H')\subset M_\pm'$, $u$ implements $\hat S$ on $H+H'$.\ev{P: assembly; S1}

\paragraph{(c) Implementer $\Rightarrow$ continuity $\Rightarrow$ Shale.} Let $UW(k)U^*=W(\hat Sk)$ on $H+H'$. As $\|W(k)\Omega-\Omega\|^2=2(1-e^{-\|k\|^2/2})$, $W(k_n)\Omega\to\Omega$ forces $k_n\to0$; conversely $k_n\to0$ gives $W(k_n)\to1$ strongly, since coherent vectors are total. So $k_n\to0$ implies $W(\hat Sk_n)=UW(k_n)U^*\to1$ and $\hat Sk_n\to0$: the real-linear $\hat S$ is continuous at $0$, hence bounded, and extends from the dense $H+H'$ to $\bar S$ on $\cK$ with $UW(k)U^*=W(\bar Sk)$. It is symplectic: $\operatorname{Im}\ip{\bar Sk}{\bar Sj}-\operatorname{Im}\ip kj\in\pi\mathbb Z$ is linear under $k\to sk$, $s$ real, hence zero. $U^*$ gives a bounded inverse. Shale's necessity direction makes $\bar\beta=\frac12(\bar S+i\bar Si)$ Hilbert--Schmidt, and $\|\bar\beta\|_{\rm HS}\ge\|\beta e\|/\|e\|$ whenever $e,ie\in H+H'$.\ev{P: assembly; ref: referee; S2}

\paragraph{(d) Sign-definite first-order floor.} Write $S(\tau)-1=gT(\tau)+R_2$, $T(\tau)=\int_0^\tau U^0(s)DU^0(-s)\,ds$, and let $\hat\beta_1$ be the antilinear part built from $T$. As $U^0$ is symplectic,
\[
\omega(k,T(\tau)k)=\int_0^\tau\omega\bigl(U^0(-s)k,DU^0(-s)k\bigr)\,ds=-m(\tau),\quad\text{where}
\]
\[
m(\tau)=\int_0^\tau\ip{\pi_s}{\GD\pi_s}ds\ \ \text{for }k=(0,\pi)\ \text{(even)},\qquad m(\tau)=\int_0^\tau\ip{q_s}{Nq_s}ds\ \ \text{for }k=(q,0)\ \text{(odd)},
\]
with $\pi_s=\cos(sh)\pi$, $q_s=\cos(sh)q$. Since $\GD\ge0$ and $N\ge0$, $m$ is odd in $\tau$ and non-decreasing for $\tau>0$: the first-order term has a definite sign and cannot vanish at isolated times.\ev{P: all-t; ref: referee} From $P+Q=Tz$ and $c_h^{-2}+s_h^{-2}=2\sinh2\pi h$, $\|\hat\beta_1k\|\ge\|(2\sinh2\pi h)^{-1/2}Tz\|\ge m(|\tau|)/(2\,\mathrm{Den})$ with $\mathrm{Den}^2=\ip z{2\sinh(2\pi h)z}$. Take $e_l=(0,a_l)$ (even) or $(a_l,0)$ (odd), $a_l=e^{-3h^2/2}h^2u$, $u(y)=(1-(\frac{y-2}4)^2)^4$ on $-2\le y\le6$, $y=\rho-\ln2\lambda$. Cauchy--Schwarz in the $\GD$- (resp.\ $N$-) form against $u$, $\cos x\ge1-x^2/2$ and Jensen's inequality on the spectral measure of $hu$ give $\lambda m(|\tau|)\ge M_{\rm lo}\Upsilon(\min(|\tau|,s_*)/s_*)$, hence $\lambda\|\hat\beta_1e_l\|/\|e_l\|\ge c(\tau)$. \emph{Uniformity in $l$:} in the variable $y$ the potential $\lambda^2/\sinh^2\rho=e^{-2y}(1-e^{-2y}/4\lambda^2)^{-2}$ and the weight $(1+\cosh\rho)/\lambda$ both decrease monotonically in $l$ to their $l=\infty$ (Liouville) limits, so every spectral quantity entering $M_{\rm lo}$, $\mathrm{Den}$ and $\|e_l\|$ lies between its values at $l=16$ and $l=\infty$. Those two endpoint values are certified by interval arithmetic, which bounds the floor for all $l\ge16$ at once.\ev{P: even-passage, odd-sign, all-t; N} The floors are not sharp: at $t=\frac12$ the grid values are about 58 (even) and 146 (odd) times larger (Table~\ref{tab:const}).

\paragraph{(e) Remainder.} The bounded perturbation theorem and Gronwall give $\|R_2\|\le e^x-1-x$, $x=\sqrt2|g||\tau|/\lambda$. With $\|\beta e-g\hat\beta_1e\|\le\frac12\|R_2\|(\|e\|+\|\coth^{3/2}(\pi h)z\|)$, a lower bound $\|e_l\|\ge E_{\rm lo}$ and an upper bound $\|\coth^{3/2}(\pi h)z_l\|\le N_3$, this gives $\|\beta e_l\|/\|e_l\|\ge c(\tau)|g|/\lambda-C_R(e^x-1-x)$ with $C_R=\frac12(1+N_3/E_{\rm lo})$, which is at least $c(\tau)|g|/(2\lambda)$ for $l\ge l_1$.\ev{P: even-remainder, odd-remainder, all-t; S2} The bounds bite only for $l$ above about $15{,}000|g|$ at $t=\frac12$; that suffices for the tail, and each is tested non-vacuously at one row.\ev{N}

\paragraph{(f) Tower sum.} The $\hat e_{l,m}=e_{l,m}/\|e_l\|$ lie in orthogonal blocks of $\cK$, so $\|\bar\beta\|_{\rm HS}^2\ge\sum_{l\ge l_1}(2l+1)(cg/2\lambda)^2=(cg/2)^2\sum_{l\ge l_1}(\frac1l+\frac1{l+1})=\infty$, contradicting (c).\ev{P: assembly}

\par\medskip\noindent\begin{minipage}{\textwidth}\centering\footnotesize
\begin{tabular}{lll}
\toprule
 & even & odd\\
\midrule
floor $c_\pi$ (second code) & $2.110\times10^{-3}$ $(2.1193\times10^{-3})$ & $2.038\times10^{-3}$ $(2.0505\times10^{-3})$\\
grid $\lambda\|\hat\beta_1e_l\|/\|e_l\|$, $t=\frac12$, $l=16\to4096$ & $0.1218\to0.1247$ & $0.2974\to0.3053$\\
$E_{\rm lo}$ (grid $\|e_l\|$), $N_3$ & $0.14972$ $(0.268)$, $0.32082$ & $0.09369$ $(0.175)$, $0.20080$\\
$C_R$ (second code) & $1.5714$ $(1.5693)$ & $1.5716$ $(1.56948)$\\
$l_1(1)$ at $t=\frac12$, formula (second code) & $14716$ $(14632)$ & $15237$ $(15124)$\\
\bottomrule
\end{tabular}
\tabcap{tab:const}{Constants of Theorem~\ref{thm:main}. $E_{\rm lo}$ and $N_3$ are certified bounds; the grid norm is shown for comparison. The second codes reproduce every first-code constant on the safe side. The direct even remainder bound at $t=\frac12$ gives the slightly smaller threshold $14714$.}
\end{minipage}\par\medskip

\paragraph{What ``normal'' excludes.} The area-side flow has a quasi-free KMS state at inverse temperature $2\pi$ on the Weyl algebra, unique among states continuous in the thermal norm, whose modular flow is that flow; by the theorem it is disjoint from the vacuum representation, and an Araki--Yamagami sum grows logarithmically. Whether it is physically acceptable (its relative entropy to the vacuum appears to diverge at the edge) is open.\ev{beyond-normal; ref: beyond-ref}

\section{The light cone}\label{sec:ruler}

On each null generator of the ball's past light cone, with affine parameter $v\in[0,R]$ from the tip and curvature datum $C(v)$, put $\sigma_p(v)=-v^{-p}\int_0^vu^pC(u)\,du$ and $\cQ_p=\frac1{4G}\oint\int_0^Rv^3(1-v/R)|\sigma_p|^2\,dv$. Given A1, the standard cone symplectic form and continuity of F1 to the boundary cone, the vacuum modular energy is $\Srel=\cQ_3$, because $2v\sigma_3$ is the unique canonical current that F1 transports locally, and the first-law area functional of the printed first-law proof is $\cQ_2$.\ev{ref: cone-ref} With $h=v^2(\sigma_2-\sigma_3)$ one has $\sigma_3=h'/v$, $\sigma_2=h'/v+h/v^2$, and one integration by parts gives
\begin{equation}\label{eq:cone}
\cQ_2-\Srel=\frac1{4G}\oint\!\int_0^Rv^3\,|\sigma_2-\sigma_3|^2\,dv .
\end{equation}
So on the cone the first-law charge exceeds the entropy for every non-zero cone curvature at every $l$, identically in both parities and on metric and extrinsic-curvature data alike; the ratio has infimum~$1$, not attained, and no upper bound.\ev{P: ruler-l; X: $l=2$--$6$} We do not claim that this settles where the asymmetry of~\eqref{eq:sector} comes from: the ingredients of the fixed-volume rule were dissected on the static tides only.

\section{Relation to the literature and discussion}\label{sec:lit}

\paragraph{Modular theory.} For the scalar the ball's vacuum modular group is the conformal flow (Hislop--Longo~\cite{HL}, a conformal image of Bisognano--Wichmann~\cite{BW}); Longo--Morsella~\cite{LM} give the modular Hamiltonian and Ciolli--Longo--Ruzzi~\cite{CLR} the coherent-state entropy. For helicity~2, Lled\'o's Prop.~5.8~\cite{Lle} shows that free nets of every helicity meet the hypotheses of Brunetti--Guido--Longo's Thm~2.3(ii)~\cite{BGL}, whose conclusion is that double-cone modular groups act as conformal flows (cyclicity by Reeh--Schlieder; the weight is not printed there; Hislop~\cite{His88} was read as an abstract only). So F1 has a published route once the curvature net is identified with Lled\'o's helicity-2 net, the smeared Weyl spinor; no paper we found names linearised gravity.\ev{Lit: primary-sources, phys-id} Benedetti--Casini~\cite{BC} use the same radial reduction and take the ball's algebra to be that of the linearised Riemann tensor; they compute entanglement entropy.

\paragraph{Jensen--Sorce--Speranza~\cite{JSS23}.} Their assumption A4 asks that the flow of a boost-like charge coincide with the modular flow of some state, a vector in their Hilbert space. On odd data at $l=2$ to $5$, their printed graviton charge, eq.~(3.28), changed under a linearised gauge transformation in our computation; their Section~3.3 with true volume fixed gave $c_lE_2$, and the covariant reading gave $c_lE_1$.\ev{X: jss-id, $l=2$--$5$} Our question to them: for non-Killing $\xi$, is the flow in A4 the transport F3 of the potential $h$, as (3.28) says, or the conformal transport F1 of the curvature? On the ball they differ by one unit of weight.

\paragraph{Satishchandran--Sorce~\cite{SS}.} Their theorem concerns the scalar, and we do not question it. Their graviton appendix says (p.~32) that their state $\omega_0^{\rm GR}$ is not the restriction of the vacuum. In $D=4$ the first four $u$-moments of the Weyl component on each cone generator vanish for every divergence-free source in the diamond (sharp; false in $D=5,6$), so their (B.26) converges, and two independent computations find it equal to the vacuum two-point function on their algebra (with the sign of their (B.13)).\ev{P: ss-lemma; X: ss-second} Is their sentence meant for $D\neq4$ only?

\paragraph{Type III, crossed products and the edge.} Jensen--Sorce--Speranza work with type II algebras obtained by adding an observer or a clock, and with states whose modular flow is geometric at an instant. Our statement concerns the free-field algebra $M$, which is of type III\ev{S3}; it excludes faithful normal states and weights of $M$, not states of a crossed-product algebra, and the quasi-free state of Section~\ref{sec:shale} lies outside the vacuum representation. Whether the obstruction survives the crossed product is open. So is its robustness: the divergence is an edge logarithm, $\sum(2l+1)/\lambda^2$, and we have not tested it against corner ambiguities of the second-order area or a smoothed edge. Measured instead as a per-sector affinity defect in a large-$l$ limit model, the obstruction has the constants $7\zeta(3)/(16\pi^2)$ (odd) and $\approx0.0135994$ (even); the proof uses neither, and two lattice analyses that agree with them were not clean blind tests.\ev{N: even-const; L: hinge-blind, even-blind} We found nothing in print that states the result; these are search results, not a priority claim.

\section{Open items}\label{sec:open}

(1)~A1: edge modes, crossed products and constraints lie outside $M$; second quantisation is cited, not proved. (2)~A2, and which charge Jensen--Sorce--Speranza mean. (3)~I, the completeness of the two data families, read in the record but not re-proved. (4)~Shale, Eckmann--Osterwalder, Engel--Nagel and Reed--Simon in the originals, and a second reading of the primary sources. (5)~Sharper floors: the bounds bite only above $l\approx15{,}000|g|$ at $t=\frac12$, each tested non-vacuously at one row. (6)~Robustness to corner ambiguities and edge smoothing; the crossed-product algebra. (7)~Dimensions $D\ge5$. (8)~No outside reader.

\paragraph{Changes since version 3.} Version~3 (10 November) compared electric with magnetic tides; version~4 replaces that contrast by metric against extrinsic-curvature data on every sector, makes the non-implementability theorem the main result, and treats version~3's ladder as its lowest rung. It withdraws version~3's three-dimensional lattice support, both the bracket $0.89$ to $1.02$ and the electric ratios $1.02$, $1.01$, $0.99$, because they came from one run whose own registered verdict was inconclusive, and keeping only the favourable half would be selective reporting. The flow on curvature, ``proved in companion notes'' in version~3, is now a refereed theorem resting on A1.

\section*{How it was checked}

Each result came from a separate lane (report, pre-registered checks with planted errors, scripts, logs); the programme's lead graded it, checked its core by hand and re-ran its scripts.

\begin{center}\footnotesize
\begin{tabular}{lp{0.38\textwidth}ll}
\toprule
Claim group & Reports & Lead's check & Second code or referee\\
\midrule
Sector identity & charge-all-l, charge-ref, even-tower & re-ran; hand & charge-ref\\
Bridge & modflow-ref, chain-all-l, chain-ref, modgen & re-ran; hand & two referees\\
Cocycle, Shale & referee, assembly, beta-def & by hand & referee\\
Floors, remainders & even-passage, odd-sign, remainders, all-t & hand; re-ran & two second codes\\
Dictionaries & even-dict, odd-dict, even-density & hand; re-ran & none\\
Light cone & ruler, ruler-l, cone-ref & hand; re-ran & cone-ref\\
Outside the folium & beyond-normal & re-ran & beyond-ref\\
Satishchandran--Sorce & ss-second, ss-tension, ss-lemma & re-ran & two routes; proof\\
Standard theorems & std-theorems, phys-id, primary-sources & identifiers & none\\
\bottomrule
\end{tabular}
\end{center}

All checkers, including the referees and the second codes, are AI models of one family working under the author's direction. Nothing here is peer reviewed.

\begin{thebibliography}{99}\footnotesize\itemsep0pt
\bibitem{JV} T.~Jacobson and M.~Visser, \emph{Gravitational thermodynamics of causal diamonds in (A)dS}, arXiv:1812.01596.
\bibitem{CLR} F.~Ciolli, R.~Longo and G.~Ruzzi, \emph{The information in a wave}, arXiv:1906.01707.
\bibitem{HL} P.~D.~Hislop and R.~Longo, Commun.\ Math.\ Phys.\ 84 (1982) 71.
\bibitem{BW} J.~J.~Bisognano and E.~H.~Wichmann, J.\ Math.\ Phys.\ 16 (1975) 985; 17 (1976) 303.
\bibitem{LM} R.~Longo and G.~Morsella, \emph{The massless modular Hamiltonian}, arXiv:2012.00565.
\bibitem{His88} P.~D.~Hislop, Ann.\ Phys.\ 185 (1988) 193.
\bibitem{Lle} F.~Lled\'o, \emph{Conformal covariance of massless free nets}, Rev.\ Math.\ Phys.\ 13 (2001) 1135, math-ph/0006018.
\bibitem{BGL} R.~Brunetti, D.~Guido and R.~Longo, Commun.\ Math.\ Phys.\ 156 (1993) 201, funct-an/9302008.
\bibitem{Der} J.~Dere\-zi\'nski, Lect.\ Notes Phys.\ 695 (2006), math-ph/0511030.
\bibitem{Shale} D.~Shale, \emph{Linear symmetries of free boson fields}, Trans.\ Amer.\ Math.\ Soc.\ 103 (1962) 149 (not opened).
\bibitem{Connes} A.~Connes, \emph{Une classification des facteurs de type III}, Ann.\ Sci.\ \'Ec.\ Norm.\ Sup.\ 6 (1973) 133.
\bibitem{Schn} R.~Schnaubelt, \emph{Evolution equations}, lecture notes, KIT (2023).
\bibitem{JSS17} T.~Jacobson, J.~M.~M.~Senovilla and A.~J.~Speranza, arXiv:1710.07379.
\bibitem{BC} V.~Benedetti and H.~Casini, \emph{Entanglement entropy of linearized gravitons in a sphere}, arXiv:1908.01800.
\bibitem{JSS23} K.~Jensen, J.~Sorce and A.~J.~Speranza, arXiv:2306.01837.
\bibitem{SS} G.~Satishchandran and J.~Sorce, arXiv:2607.20678 (July 2026).
\bibitem{T} A.~Houg, the programme's record: lane reports, scripts and logs, waves 64 to 73 (November 2026), \url{https://github.com/ni1ra/graviton-ladder}, folders under \texttt{record/wave73/}. By row of Table~\ref{tab:hyp}: proved: charge-all-l, charge-ref, even-dict, odd-dict, even-proof, hinge-proof, even-density, referee, beta-def, assembly, even-passage, odd-sign, even-remainder, odd-remainder, all-t, even-second, odd-second; bridge: wave70 modflow, modflow-ref, chain-all-l, chain-ref, modgen, subspace; standard: primary-sources, std-theorems, phys-id.
\end{thebibliography}

\end{document}
